\subsubsection{BlackBerry OS}

\paragraph{Resources Managed}

BlackBerry~10 manages the CPU, memory, I/O devices, and files. A
distinctive design point is that device drivers and file systems are
themselves ``resource managers''---user-space processes that present
devices and files through the unified pathname space
\citep{qnx_memman, qnx_powersafe_fs}. Resources fall into the usual
categories: \emph{sharable} (e.g. read-only code, the CPU over time)
versus \emph{non-sharable} (e.g. a device being written), and
\emph{preemptable} (CPU, memory) versus \emph{non-preemptable} (a
partially completed device transaction). In a real-time phone, CPU
time and responsiveness are the most critical resources, which is why
so much of QNX's resource machinery is about guaranteeing CPU shares
\citep{qnx_aps}.

\paragraph{Resource Allocation Strategies}

Allocation is dynamic and priority-driven by default---the
highest-priority \textsc{ready} thread gets the CPU, memory is
allocated on demand, and resource managers serve clients in priority
order \citep{qnx_kernel_sysarch, qnx_memman}. Layered on top is an optional
adaptive partitioning scheduler (APS), which lets the designer reserve
a guaranteed minimum percentage of CPU (a budget) for a group of
threads/processes \citep{qnx_aps}. APS is a configurable QNX
feature---widely used in domains such as automotive---rather than a
confirmed default of BlackBerry~10; the discussion below therefore
describes a capability the platform provides and a BlackBerry~10 build
could enable, not a documented handset setting. Critically, partitions
are \emph{adaptive}: when the system is under-loaded, a partition can
borrow idle CPU beyond its budget; only when the system is overloaded
are the guaranteed minimums strictly enforced
\citep{qnx_aps}. This combines high utilisation (no CPU
wasted) with hard guarantees (important work always gets its share).

\paragraph{Resource Tracking}

The kernel tracks thread/process state and CPU consumption; when
enabled, APS keeps sliding-window accounting of each partition's CPU
usage (window 8--400\,ms, plus longer statistics windows viewable via
the \texttt{aps} command) to enforce budgets
\citep{qnx_aps, qnx_aps_cpu_percent}. Memory is tracked through the
process manager's page tables, and resources are reclaimed
automatically when a process terminates---its address space,
descriptors, and connections are released \citep{qnx_memman}. A neat
consequence of Send/Receive/Reply is that when a resource manager
(driver/file system) does work on behalf of a client, the CPU time
would---under APS---be \emph{billed to the client's partition}, so
tracking reflects who actually caused the load \citep{qnx_aps}. Tracking adds modest
overhead, kept bounded by the windowed averaging scheme
\citep{qnx_aps_cpu_percent}.

\paragraph{Load Balancing and Distribution}

On multicore hardware, BlackBerry~10 distributes threads across cores
under symmetric multiprocessing (SMP), with processor-affinity controls
that let the designer bind threads to specific cores to balance cache
locality against spreading load \citep{qnx_smp_multicore}. Where APS
is deployed, it adds \emph{temporal} load balancing---distributing CPU
time fairly across partitions over its averaging window so that no
function is starved while another hogs the processor \citep{qnx_aps}. The performance
implication is high, controlled utilisation; the risk it manages is
exactly the ``some resources overused while others idle'' problem the
guide raises.

\paragraph{Handling Deadlock}

QNX's primary stance is avoidance/prevention by design rather than a
Banker's-algorithm detector. The synchronous Send/Receive/Reply model
imposes a disciplined, predictable communication ordering that reduces
the conditions for deadlock \citep{qnx_ipc}. Priority inheritance on
mutexes and messages prevents priority inversion (a deadlock-adjacent
hazard) by temporarily boosting a blocking holder's priority
\citep{qnx_kernel_sysarch, qnx_ipc}. For situations that still go
wrong---including a hung or deadlocked server---QNX's architecture
allows the offending process to be restarted without taking down the
system, and its optional HA Framework provides a watchdog (the HAM)
that can detect unresponsive processes and restart them as a
system-level backstop where deployed \citep{qnx_ham}. The trade-off: QNX does not centrally prevent
all deadlocks algorithmically (that burden sits partly with the
application designer), but its architecture makes deadlock both less
likely and recoverable.

\paragraph{Resource Contention and Resolution}

When processes compete for a resource, requests are served in priority
order, and a mutex is handed to the highest-priority waiting thread on
release, so threads pass through a critical region in priority order
\citep{qnx_kernel_sysarch}. Priority inheritance keeps a high-priority
requester from being blocked indefinitely behind a low-priority holder
\citep{qnx_ipc}. Under high contention, an APS-enabled QNX configuration guarantees
each partition its floor so contention cannot completely lock out
important work \citep{qnx_aps}. Fairness is thus \emph{priority-weighted} rather than
egalitarian---appropriate for a real-time system.

\paragraph{Performance and Efficiency}

Resource management adds overhead in two places: message-passing and
context-switch cost for every resource-manager interaction
\citep{qnx_arm_memory}, and the windowed accounting of APS
\citep{qnx_aps_cpu_percent}. QNX balances efficiency against fairness
by letting partitions borrow idle time (efficiency) while enforcing
minimum budgets only under load (fairness), and balances performance
against complexity by making partitioning \emph{optional}---designers
pay for it only when they need the guarantees \citep{qnx_aps}.

\paragraph{Scalability}

BlackBerry~10 scales with added processes/users bounded mainly by
physical resources, and scales across cores via SMP with processor
affinity \citep{qnx_smp_multicore}. APS, where deployed, scales to many
partitions and has been used in complex systems (e.g. automotive) where
dozens of subsystems share a chip \citep{qnx_aps}. As
load rises, the limits that appear are message-passing overhead and
the per-process memory cost of private address spaces
\citep{qnx_memman, qnx_arm_memory}---the price of the isolation that
makes the system robust.

\paragraph{Overview --- Resource Management}

Strengths include dynamic priority allocation, optional
\emph{guaranteed} CPU budgets via adaptive partitioning where enabled,
automatic reclamation, accurate client-billed accounting, and
architectural restartability (optionally managed by the HA Framework's
HAM) \citep{qnx_memman, qnx_ham, qnx_aps}. Weaknesses are
message-passing/accounting overhead and reliance on the designer to
structure partitions and avoid application-level deadlock
\citep{qnx_aps, qnx_arm_memory}. Under heavy load, a QNX system using
APS behaves predictably: low-priority/untrusted partitions are
throttled to their minimums so critical work keeps running, whereas a
purely priority-based competitor could let high-priority work starve
everything else \citep{qnx_aps}. In a resource-constrained system,
QNX's small footprint (an HMI can run in 32\,MB or less) and explicit
budgeting make it well suited, often better than a heavier
general-purpose OS \citep{slideshare_rim_qnx}.
